\documentclass{ximera}

% \input{../preamble.tex}

%\title{Constant and Linear Functions}

\begin{document}
	\author{Stitz-Zeager}
	\xmtitle{Constant and Linear Functions}
\license{CC BY-NC-SA 4.0}   
\begin{abstract}
\end{abstract}
%\maketitle

%\mfpicnumber{1}
%\opengraphsfile{ConstantAndLinearFunctions}





 

\subsection{Linear Regression}
\label{LinearRegression}
 
We have demonstrated in this section that constant, linear, and piecewise combinations of these two function types can be used to model a variety of phenomena inspired by real-world situations.  What happens, as if often the case in real-world situations, when we are given data sets that are not precisely linear, but still have a definite linear trend?  An example of this is Skippy's  time and temperature data from Example \ref{timetempex1} in Section \ref{FunctionsandtheirRepresentations}. 



In that example, $t$ represented the time (number of hours after 6 a.m.) and $T$ represented the outdoor temperature in degrees Fahrenheit. The data Skippy collected along with a plot of the function $T = f(t)$ are given below. Even though the data points as $t$ varies from $t = 0$ to $t = 8$ do not all lie on the same line - a fact we could prove analytically by checking slopes - there does appear to be a linear \textbf{trend} evident.  The same can be said for the data as $t$ varies from $t =8 $ to $t = 12$.  As we'll see, there are statistical methods which can produce  linear functions that are in some sense `closest' to all of the data, and they are represented below by the dashed lines below on the right. 

\begin{center}
$\begin{array}{|c||c|}  \hline

 \text{$t$: hours after 6 a.m.}  & \text{$T$: temperature $^{\circ}$F} \\ \hline
 0 & 64  \\
 2 & 67  \\
 4 &  75  \\
 6 &  80 \\
 8 & 83  \\
 10 &  83 \\
 12 & 82  \\
\end{array}$
\end{center}
% https://www.desmos.com/calculator/aqedfuqmuw

\begin{center}
    \desmos{aqedfuqmuw}{800}{600}
\end{center}


How do we measure how `close' a set of points is to a given line?   Let's leave Skippy's data for the moment and focus on a smaller data set.  Suppose we collected three data points: $\{(1, 0.5), (3, 2), (4, 3)\}$.  At the top of the next page (on the left) we plot these points along with the line $y = 0.5 x + 0.5$.   The way we measure how close the line is to these points is by computing the \index{regression ! total squared error}\index{total squared error}\textbf{total squared (vertical) error} between the data points and the line as follows.   For each of our data points, we find the vertical distance between the point and the line.  To accomplish this, we need to find a point on the line directly above or below each data point.  In other words, we need a point on the line with the same $x$-coordinate as our data point.  



For example, to find the point on the line directly above $(1,0.5)$, we plug $x=1$ into $y=0.5x + 0.5$ and we get the point $(1,1)$.  Similarly, we find $(3,2)$ is on the line already and $(4,2.5)$ is the point on the line directly beneath $(4,3)$.   We find the total squared error $E$ by taking the sum of the squares of the differences of the $y$-coordinates of each data point and its corresponding point on the line. For the data and line in this discussion $E = (0.5-1)^2+(2-2)^2+\left(3-2.5\right)^2 = 0.5$.  



Using advanced mathematical machinery,\footnote{like Calculus or Linear Algebra \ldots} it is possible to find the line which results in the lowest value of $E$.  This line is called the \index{regression ! least squares line}\index{least squares regression line}\index{line ! least squares regression}\textbf{least squares regression line}, or sometimes the `line of best fit'. \index{line ! of best fit} The formula for the line of best fit requires notation we won't present until Chapter \ref{SequencesandtheBinomialTheorem}, so we will revisit it then.  Most graphing utilities have a built-in regression feature, so at this point we turn the computations over to the technology.  A screenshot from \href{www.desmos.com}{\underline{desmos}} is given on the right at the top of the next page.  





\begin{center}
    \begin{tikzpicture}
\begin{axis}[
    axis lines=middle,
    xmin=-1, xmax=5,
    ymin=-1, ymax=5,
    xtick={1,2,3,4},
    ytick={1,2,3,4},
    xlabel={$x$},
    ylabel={$y$},
    axis equal image,
    every axis x label/.style={at={(ticklabel* cs:1)},anchor=west},
    every axis y label/.style={at={(ticklabel* cs:1)},anchor=south},
    width=8cm, height=8cm,
    enlargelimits=false,
]

% The function y = 0.5x + 0.5
\addplot[domain=-1:5, thick] {0.5*x+0.5};

% Dotted vertical segments
\addplot[dotted, thick] coordinates {(1,0.5) (1,1)};
\addplot[dotted, thick] coordinates {(4,3) (4,2.5)};

% Points
\addplot[only marks, mark=*] coordinates {
    (1,0.5) (1,1) (3,2) (4,3) (4,2.5)
};

\end{axis}
\end{tikzpicture}

% https://www.desmos.com/calculator/ofrcaiemar

\desmos{ofrcaiemar}{800}{600}
\end{center}

Our graphing utility produces the model\footnote{We chose to use three decimal places for the approximations in this demonstration.  How many you get to use in reality varies from one application to another.} $y=mx+b$ where the slope is $m \approx 0.821$ and the $y$-coordinate of the $y$-intercept is $b \approx -0.357$.  The value $r$ is the \index{regression ! correlation coefficient}\index{correlation coefficient}\textbf{correlation coefficient} and is a measure of how close the data is to being on the same line.  The closer $|r|$ is to $1$, the better the linear fit.\footnote{The value $r^2$ is called the \index{regression ! coefficient of determination}\index{coefficient of determination}\textbf{coefficient of determination} and is also a measure of the goodness of fit. We refer the interested reader to a course in Statistics to explore the significance of $r$ and $r^2$.} Having $r \approx 0.997$ tells us that the points have a strong, positive correlation - that is, they are very close to being on a line with a positive slope, namely $y = 0.821x - 0.357$.   Indeed, the total squared error between our data set and this line is $E \approx 0.018$. The mathematics tells us that this is the smallest we can get $E$ by modifying the parameters $m$ and $b$, even though none of the data points actually lie on the line.  



Now that we have this new mathematical machinery, let's revisit Skippy's time and temperature data.
  
\begin{example} \label{timetempregressionex}  $~$
  
\begin{enumerate}
  
\item Use a graphing utility to find best fit linear models for each of the data sets below. Comment on the fit and interpret the slope of each.

\begin{center}
  

  
$\begin{array}{|c||c|}  \hline

 \text{$t$: hours after 6 a.m.}  & \text{$T$: temperature $^{\circ}$F} \\ \hline
 0 & 64  \\
 2 & 67  \\
 4 &  75  \\
 6 &  80 \\
 8 & 83  \\

\end{array}$\quad\quad
$\begin{array}{|c||c|}  \hline

 \text{$t$: hours after 6 a.m.}  & \text{$T$: temperature $^{\circ}$F} \\ \hline

 8 & 83  \\
 10 &  83 \\
 12 & 82  \\

\end{array}$ 
  

  
\end{center}

\item  Use your models to predict the temperature at $7$ a.m. and $3$ p.m., rounded to one decimal place.
  
\end{enumerate}

\begin{explanation}
  
\begin{enumerate}
  
 \item For our first set of data, we get the line $T = F(t) = 2.55t + 63.6$.  The value $r = 0.987$ tells us that it is a fairly good fit and we see this graphically, too.\footnote{We use $F$ as the name of the function here to distinguish it from $f$ - the function determined solely by the given set of data.}   Thus we can be confident in using this model to predict the temperature during between the hours of 6 a.m. and 2 p.m. with reasonable accuracy.   



To interpret the slope, we recognize $t$ as the independent variable (input) and $T$ as the dependent variable (output), so the slope $m = \frac{\Delta T}{\Delta t}$ is the rate of change of temperature with respect to time.   In this case,  $m = 2.55$ means that the temperature is increasing (getting warmer) at a rate of $2.55^{\circ}$F per hour.  A screenshot from Desmos is given below.

% https://www.desmos.com/calculator/jhkc4szlj7
\begin{center}
\desmos{jhkc4szlj7}{800}{600}
\end{center}
  
For the second set of data, we get $T = G(t) = -0.25t + 85.167$ and we have $r = -0.866$.  Here, the negative sign on $r$ indicates a negative correlation which means our line has a negative slope.\footnote{We use $G$ as the function name here to distinguish it from the given function $f$ and the regression for the first data set, $F$.} While the fit looks OK, it certainly isn't as strong as with the first data set, so using this model to predict the temperature between $2$ p.m. and $6$ p.m. (let alone beyond) is a bit risky.  



The slope in this case is $m = -0.25$ which corresponds to the temperature decreasing (getting cooler) at a rate of  $0.25^{\circ}$F per hour.  That's what a negative correlation means - an increase in input (more time passes) yields a decrease in output (cooler temperatures).  As screenshot from Desmos is given at the top of the next page.

% https://www.desmos.com/calculator/rqhhlc0enf  
\begin{center}
\desmos{rqhhlc0enf}{800}{600}
\end{center}

\item  The time $7$ a.m. corresponds to $t = 1$. This falls between $t = 0$ and $t = 8$ so we use our first model.  Substituting $t = 1$ gives $T = F(1) = 2.55(1) + 63.6 \approx 66.2$.  Therefore, the model predicts the temperature to be $66.2^{\circ}$F at 7 a.m..  Likewise, $3$ p.m. corresponds to $t = 9$.  This is greater than $8$, so we use the second model: $T = G(9) =  -0.25(9) + 85.167 \approx  82.9$.  The model predicts the temperature at $3$ p.m. to be $82.9^{\circ}$F.   Based on the goodness of fit of each model, we have more confidence in the former prediction than in the latter.   
  
\end{enumerate}
\end{explanation}
\end{example}
 
Examples \ref{PortaBoyCost}, \ref{PortaBoyDemand} and \ref{timetempregressionex} (among others)  represent three different levels of mathematical modeling.  In Example  \ref{PortaBoyCost},  the mathematical model (the cost function) was provided and our task was to use the model to \textbf{interpret} the mathematics in that context.    In Example  \ref{PortaBoyDemand}, we were given a minimal amount of information, namely, two data points, and then asked to \textbf{construct} a model which fit those data exactly.  Lastly, in Example \ref{timetempregressionex}, we were given several data points and we used statistical methods to construct a \textbf{best fit} model to the data. 



The validity of the models rests on the validity of the underlying assumptions used to create the models.  For instance, is there any reason to assume a price-demand function would be linear?  Is it reasonable to assume that the temperature changes at a constant rate? These are questions for economists and scientists.  Mathematicians often take on a role of equal parts translator and prophet:  they codify ideas into formulas and then use them to make predictions about yet-to-be observed phenomena.  


\subsection{The Average Rate of Change of a Function}
\label{AverageRateofChange}

As mentioned earlier in the section, the concepts of slope and the more general rates of change are important concepts not just in Mathematics, but also in other fields.   Many important phenomena are modeled using non-linear functions, and while the rates of change of these functions are not constant, we can sample the function at two points and compute what is known as an \textbf{average rate of change} between them to give some sense as to the function's behavior over that interval.\footnote{We are basically pretending that the function is linear on a short interval to see what we can say about its behavior.}





\begin{definition} \label{arc}  Let $f$ be a function defined on the interval $[a,b]$. 

 The \index{average rate of change}\index{rate of change ! average}\textbf{average rate of change}  of $f$ over $[a,b]$ is defined as: \[ \frac{\Delta [f(x)]}{\Delta x} = \frac{f(b) - f(a)}{b-a} \]

Geometrically, the average rate of change is the slope of the line\footnote{This line is called a \index{secant line}\index{line ! secant}\textbf{secant line}.}  containing $(a, f(a))$ and $(b, f(b))$.

\end{definition}

% https://www.desmos.com/calculator/2ctplusi5g

\begin{center}
\desmos{2ctplusi5g}{800}{600}
\end{center}

As with Definitions \ref{absmaxmindefn}  and \ref{incdeccnstdefn}, the wording in Definition \ref{arc}, while referring to the function $f$, is really making a statement about its outputs $f(x)$.  



If $f$ is increasing over $[a,b]$, then the average rate of change will be positive. Likewise, if $f$ is decreasing or constant, the average rate of change will be negative or $0$, respectively. (Think about this for a moment.) However, as the next example demonstrates, the converses of these statements aren't always true.\footnote{For example, the average rate of change over an interval could be positive yet the function could decrease over part of that interval and then increase on a different part.}

\begin{example} \label{ARCRocketExample} The formula $s(t) = -5t^{\,2} + 100t$ for $0 \leq t \leq 20$ gives the height, $s(t)$, measured in feet, of a model rocket above the Moon's surface as a function of the time $t$, in seconds after lift-off.

\begin{enumerate}

\item Find $s(0)$, $s(5)$, $s(10)$, $s(15)$ and $s(20)$ and use these along with a graphing utility to graph $y = s(t)$.

\item State the range of $s$ and interpret the extrema, if any exist.

\item Find and interpret the $t$- and $y$-intercepts.  

\item Find and interpret the interval(s) over which $s$ is increasing, decreasing or constant.  

\item Find and interpret the average rate of change of $s$ over the intervals $[0,5]$, $[5,10]$, $[10, 20]$ and $[5, 15]$.


\end{enumerate}

\begin{explanation}

\begin{enumerate}

\item  To find $s(0)$, we substitute $t=0$ into the formula for $s(t)$:  $s(0) = -5(0)^2+100(0) = 0$.  Similarly,  $s(5) = -5(5)^2+100(5) = -5(25)+500 = -125+500 = 375$.  Continuing, we obtain: $s(10) = 500$, $s(15) = 375$ and $s(20) = 0$.  We construct a table of values and with a graphing utility we obtain:

% https://www.desmos.com/calculator/gh2t8utkoz

\begin{center}
\desmos{gh2t8utkoz}{800}{600}
\end{center}

\item  Projecting the graph to the $y$-axis, we see that the range of $s$ is $[0, 500]$ so the minimum of $s$ is $0$ and the maximum is $500$.  This means that the rocket at some point is on the surface of the Moon and reaches its highest altitude of $500$ feet above the lunar surface.

\item The first intercept we see is $(0, 0)$ which is both a $t$- and a $y$-intercept.  Since $t$ is the time after lift-off and $y = s(t)$ is  the height above the Moon's surface, the point $(0, 0)$ means that the model rocket was launched ($t=0$) from the Moon's surface ($s(t) = 0$).  The remaining intercept, $(20, 0)$,  is another $t$-intercept.  This means that $20$ seconds after lift-off ($t=20$), the model rocket returns to the Moon's surface ($s(t) = 0$).   That is, the `time of flight' of the model rocket is 20 seconds.

\item  Referring to Definition \ref{incdeccnstdefn}, $s$ increases over the interval $[0, 10]$, since for those values of $t$, as we read from left to right, the graph of the function is rising meaning the $y$ values (hence $s(t)$ values) are getting larger. Thus the model rocket is heading upwards for the first $10$ seconds of its flight.  We find that $s$ decreases over the interval $[10, 20]$, indicating once it has reached its highest altitude of $500$ feet $10$ seconds into the flight, the rocket begins to fall back to the surface of the Moon, landing $20$ seconds after lift-off.

\item  To find the average rate of change of $s$ over the interval $[0, 5]$ we compute \[ \frac{\Delta[s(t)]}{\Delta t} = \frac{s(5) - s(0)}{5 - 0} = \frac{\text{$375$ feet}}{\text{$5$ seconds}} = \text{$75$ feet per second}.\] In other words, the height is \textbf{increasing} at an \textbf{average rate} of $75$ feet per second during the first 5 seconds of flight.  The rate here is called the\index{average velocity} \index{velocity ! average} \textbf{average velocity} of the rocket over this interval.  Velocity differs from speed in that velocity comes with a direction.  In this case, a positive velocity indicates that the rocket is traveling \textbf{upwards}, since when $s$ is increasing, the model rocket is climbing higher.  



Similarly, the average rate of change of $s$ over the interval $[5, 10]$ works out to be $25$.  This means that the average velocity over the next $5$ seconds of the flight has slowed to $25$ feet per second. The model rocket is still, on average, traveling upwards, albeit more slowly than before.  



Over the interval $[10, 20]$, the average rate of change of $s$ works out to be $-50$.  This means that, on average, the rocket is \textbf{falling} at a rate of $50$ feet per second.  The rocket has managed to fall from its highest point $500$ feet above the surface of the Moon back to the Moon's surface in $10$ seconds so this makes sense. Finally, the average rate of change of $s$ over $[5, 15]$ is $0$.  This means that the model is the same height above the ground after $5$ seconds ($375$ feet) as it is after $15$ seconds.    



Geometrically, the average rate of change of a function over an interval can be interpreted as the slope of a secant line. Below on the left is a dotted line containing $(0, 0)$ and $(5, 375)$ (which has slope $75$) along with a dotted line containing the points $(5, 375)$ and $(10, 500)$ (which has slope $25$).  Visually, the lines help demonstrate that, while $s$ is increasing over $[0, 10]$, the rate of increase is slowing down as $t$ nears $10$.  

% https://www.desmos.com/calculator/qj3xoaviev
\begin{center}
    \desmos{qj3xoaviev}{800}{600}
\end{center}





The graph above on the right depicts a dotted line through $(10, 500)$ and $(20,0)$ indicating a net decrease over that interval.  We also have a horizontal line ($0$ slope) containing the points $(5, 375)$ and $(15, 375)$, which shows no net change between those two points, despite the fact that the rocket rose to its maximum height then began its descent during the interval $[5, 15]$.   




\end{enumerate}

\end{explanation}

\end{example}

An important lesson from the last example is that average rates of change give us a snapshot of what is happening \textbf{at the endpoints} of an interval, but not necessarily what happens \textbf{over the course} of the interval.  Calculus gives us tools to compute slopes \textbf{at} points which correspond to  \textbf{instantaneous} rates of changes.  While we don't quite have the machinery to properly express these ideas, we can hint at them in the Exercises. Speaking of exercises \ldots 

\end{document}
